\documentclass[10pt,a4paper]{article}
\usepackage[utf8]{inputenc}
\usepackage[a4paper,left=3cm,right=2cm]{geometry}
\usepackage{amsmath}
\usepackage[thmmarks, amsmath, thref]{ntheorem}
\usepackage{amsfonts}
\usepackage{amssymb}
\usepackage{fancyhdr}
\pagestyle{fancy}
\usepackage{anysize}
\usepackage{graphicx}
\usepackage{float}
\usepackage{hyperref}
\usepackage{enumitem}
\newtheorem{theorem}{Theorem}
\theoremsymbol{\ensuremath{\blacksquare}}
\newtheorem*{solution}{Solución:}

\usepackage{xcolor}

% Margenes y formalidades

\textwidth 6.25in % Margen Derecho y ancho del texto
\textheight 8.5in
\oddsidemargin 0in % Margen Izquierdo
\headheight 0in % Largo del texto
\leftmargin 1in
\parindent 0em % Salto entre la indentación (Sangría)https://www.overleaf.com/4519127669sdpptkhvsbfn
\topmargin -0.5in % Margen Superior
\parskip 0.5ex % Salto entre parrafos
\baselineskip 1pt



\lhead{
\includegraphics[scale=0.1]{figuras/logoUSM.png}} % texto izquierda de la cabecera
%\chead{}                                      % texto centro de la cabecera
\rhead{Universidad Técnica Federico Santa María\\Ingeniería Comercial } % número de página a la derecha
\renewcommand{\headrulewidth}{0.4pt} % grosor de la línea de la cabecera
\renewcommand{\footrulewidth}{0.4pt} % grosor de la línea del pie

\begin{document}
\vspace{5cm}

\begin{center}
    {\Large  Ejercicios 5: Econometría }

Profesor:    Pedro Comte Hidalgo
\end{center}



\section{Ejercicios}
\begin{enumerate}
    \item Supongamos un modelo de Regresión Lineal Simple, que describe el salario en función de la educación de la forma:
    $$salary = \beta_0 + \beta_1 educ$$
    Si la solución mediante MCO, los entrega la siguiente forma:
        $$\hat{salary} = 10 + 2\cdot educ + \hat{u}$$
    siendo $\hat{u}$ el error observado.
    \begin{enumerate}
        \item  ¿Como se interpreta el valor de $\hat{\beta_1}$?  \\
        

        \item Si dada la gran escala de $salary$, se le aplica $log()$, para tener valores más faciles de trabajar, encontrando la siguente RLS mediante MCO:
        $$\hat{log({salary})} = 2.8 + 0.07\cdot educ + \hat{u}$$
        ¿Como se interpreta el valor de $\hat{\beta_1}$ en este caso?  \\
        
        \item Ahora, se le aplica $log()$ también a educ, para tener valores más faciles de trabajar, encontrando la siguente RLS mediante MCO:
        $$\hat{log({salary})} = 1.4 + 0.80\cdot log(educ) + \hat{u}$$
        ¿Como se interpreta el valor de $\hat{\beta_1}$ en este caso? ¿En economia, como se le llama a este valor?\\
        
    \end{enumerate}


\item Se tiene un modelo que explica el puntaje en la PAES dado el IQ del individuo, de la sigueinte forma:
$$Puntaje = 200 + 10 \cdot IQ -  0.05IQ^2$$
\begin{enumerate}
    \item ¿Cual es el efecto marginal del IQ en el puntaje?
    
    \item ¿Cual es el valor donde el efecto marginal del IQ es cero?
    
    \item ¿Que sucede para valores de IQ mayores a 100?
    
\end{enumerate}

\item Se tiene la siguiente RLS:
$$y = \beta_0 + \beta_1x$$
mediante MCO podemos encontrar $\hat{\beta_0}$ y $\hat{\beta_1}$ insesgados y eficientes. Si reescalamos $\Tilde{x} = 10x$. 
\begin{enumerate}
    \item Demuestre que $\hat{\beta_1}=10\hat{\Tilde{\beta_1}}$. \\
    
\item Demuestra que el valor $t$ del estimador no cambia. \\

    
\end{enumerate}

\item Se tiene los siguentes resultados de un modelo de Regresión Lineal, de las notas finales de un curso, dada caracteristicas del alumno. De los 674 datos se tienen los siguientes resultados de los estimadores:

\begin{center}
    \includegraphics[scale=1]{Ejercicios/sem1-2/eje5.jpg}
\end{center}

\begin{enumerate}
\item Con un nivel de significancia del 5\%, defina cuales son las variables que son significativamente distintas de cero.
\\

\item Haga el intervalo de confianza para el estimador de la variable $attend$, al 95\% de confianza. Concluya.\\
\item Haga el intervalo de confianza para el estimador de la variable $termGPA$, al 95\% de confianza. Concluya.\\
\end{enumerate}

\item Se tiene los siguientes resultados de un modelo de 5 variables:
\begin{center}
    \includegraphics[scale=0.7]{Ejercicios/sem1-2/eje5.2.jpg}
\end{center}
Luego se sacan dos variables y se obtienen los siguientes resultados:
\begin{center}
    \includegraphics[scale=0.7]{Ejercicios/sem1-2/eje5.3.jpg}
\end{center}
Se quiere saber si las variables excluidas del modelo son relevantes para el mismo, con una significancia del 5\% con valor critico de $F_{95\%,2,668}=3.01$, siendo 2 las variables que se eliminan, y n-k-1=668



\end{enumerate}

    
\end{document}